\documentclass[11pt,a4paper]{article}

\usepackage[utf8]{inputenc}
\usepackage{amsmath}
\usepackage{amssymb}
\usepackage[bookmarks=false]{hyperref}
\usepackage{booktabs}
\usepackage{listings}
\usepackage{geometry}
\usepackage{url}
\geometry{margin=1in}
\usepackage{xcolor}

\definecolor{codegreen}{rgb}{0,0.6,0}
\definecolor{codegray}{rgb}{0.5,0.5,0.5}
\definecolor{codepurple}{rgb}{0.58,0,0.82}
\definecolor{backcolour}{rgb}{0.95,0.95,0.92}

\lstdefinestyle{mystyle}{
    backgroundcolor=\color{backcolour},
    commentstyle=\color{codegreen},
    keywordstyle=\color{magenta},
    numberstyle=\tiny\color{codegray},
    stringstyle=\color{codepurple},
    basicstyle=\ttfamily\small,
    breakatwhitespace=false,
    breaklines=true,
    captionpos=b,
    keepspaces=true,
    numbers=none,
    showspaces=false,
    showstringspaces=false,
    showtabs=false,
    tabsize=4
}
\lstset{style=mystyle}

\title{\textbf{liteLRU: A Cache-Coherent, Lock-Free Approximate LRU via\\
Chunked Bitmask Eviction on Modern Multi-Core Hardware}}
\author{
  xDarkicex \\
  \small Affiliations: libravdb $\cdot$ bitdev $\cdot$ zephyr-systems \\
  \small Contact: \href{mailto:git@libravdb.com}{git@libravdb.com} \\
  \small Repository: \url{https://github.com/xDarkicex/liteLRU}
}
\date{}

\begin{document}

\maketitle

\begin{abstract}
We describe the design and implementation of \texttt{liteLRU}, a fixed-capacity approximate LRU cache targeting multi-core concurrent workloads. The fundamental design question is not \textit{which items to evict}, but \textit{how to perform eviction and recency tracking without serializing the CPU cores that share the cache}. We analyze why conventional data structures---doubly-linked lists, per-node reference counting, and Array of Structs memory layouts---are structurally incompatible with the MESI cache coherency protocol~\cite{mesi1984} under concurrent access. We justify the selection of off-heap memory via mmap-backed allocators to eliminate garbage collection interference with tail latency~\cite{gogc2018}. We derive a Structure of Arrays (SoA) layout from first-principles cache-line occupancy arguments and derive the Chunked Bitmask CLOCK eviction algorithm~\cite{clock1969} from the observation that tracking recency state across 64 slots fits exactly into a 64-bit integer, enabling $\mathcal{O}(1)$ bulk state evaluation using a single hardware \texttt{CTZ} instruction~\cite{inteltzcnt,armclz}. We contrast our approach against recent eviction policy research including S3-FIFO~\cite{s3fifo2023} and SIEVE~\cite{sieve2024}, noting that policy-level optimizations are orthogonal to and do not resolve the synchronization bottleneck addressed here. Empirical evaluation on an 8-core ARM architecture yields a p99.9 latency of 1.4~$\mu$s under a contended parallel workload.
\end{abstract}

\section{Background: The Hardware Model}

To reason about concurrent cache design, we first establish a precise model of the hardware execution environment.

\subsection{The CPU Cache Hierarchy}

A modern multi-core processor maintains a private L1 cache (typically 32--64 KiB, 4--5 cycle latency), a private L2 cache (256 KiB -- 1 MiB, $\sim$12 cycle latency), and a shared L3 cache (several MiB, $\sim$40 cycle latency). Main memory sits at approximately 200--300 cycles. The performance implication is stark: a single L1 cache miss costs roughly $40\times$ more than an L1 hit.

The fundamental unit of coherency is the \textbf{cache line}, universally 64 bytes on x86\_64 and ARM64. The CPU never fetches individual bytes; it fetches and writes back entire 64-byte lines.

\subsection{The MESI Coherency Protocol}

The MESI protocol (Modified, Exclusive, Shared, Invalid) was introduced by Papamarcos and Patel~\cite{mesi1984} to maintain coherency across private caches in a multi-processor system. The relevant state transitions are:

\begin{itemize}
    \item A line held in \textbf{Shared} state by multiple cores transitions to \textbf{Invalid} in all other cores when any single core writes to it.
    \item A write to an \textbf{Invalid} line requires a Read For Ownership (RFO)---a broadcast on the CPU interconnect requesting exclusive ownership from all other cores.
\end{itemize}

The critical consequence is \textbf{false sharing}~\cite{falsesharing1993}: if two independent variables, accessed concurrently by two different cores, reside on the same 64-byte cache line, a write by one core forces an RFO on the other core's copy, even though the second core is accessing a logically unrelated variable. This is a pure hardware serialization artifact with no algorithmic solution other than spatial separation.

\subsection{Atomics and Memory Ordering Cost}

Atomic operations (\texttt{CAS}, fetch-and-add) on a memory location whose cache line is not in the M (Modified) state require:
\begin{enumerate}
    \item An MESI state transition to Exclusive via RFO.
    \item The atomic read-modify-write in the modified cache line.
\end{enumerate}

Under high concurrency, a single globally contested atomic (e.g., a global clock-hand pointer) causes every participating core to serialize on RFO round-trips. As core count $N$ grows, the wait time for the contested line scales as $\mathcal{O}(N)$ per operation, a fact formalized by Amdahl~\cite{amdahl1967}. The lock-free design methodology necessary to escape this bound is treated extensively by Herlihy and Shavit~\cite{lockfree2004}.

\section{Why Off-Heap Memory}

Go's garbage collector (GC) is a concurrent tri-color mark-and-sweep~\cite{gogc2018}. It introduces two categories of tail-latency interference directly relevant to a cache implementation:

\begin{itemize}
    \item \textbf{Write Barrier Overhead.} The GC requires write barriers on every pointer store to maintain the tri-color invariant. For a structure that stores handler function pointers or string headers, every \texttt{Add} operation incurs GC bookkeeping cost proportional to the number of pointer-typed fields written.
    \item \textbf{Stop-the-World Pauses.} While modern Go GC pause times are sub-millisecond, they are non-zero and non-deterministic. For a component targeting p99.9 tail latency in the low-microsecond range, any unpredictable pause invalidates latency guarantees.
    \item \textbf{Mark-Phase Heap Pressure.} A cache holding thousands of string values generates a large root set for the GC to scan during the mark phase, increasing mark duration proportionally with capacity.
\end{itemize}

\texttt{liteLRU} addresses this through two mechanisms:

\begin{enumerate}
    \item \textbf{Off-Heap HashMap:} An open-addressed hash map backed by an mmap-allocated region~\cite{mmap2015}. Because mmap allocations are outside the Go heap, the GC scanner does not traverse them, eliminating pointer scanning overhead for the entire index structure.
    \item \textbf{Fixed-Width Value Arrays:} The SoA slices holding methods, paths, and params are allocated once at initialization time. They are fixed-assignment slots rather than append-growing structures, bounding GC root growth to a constant determined at construction.
\end{enumerate}

\section{Why Structure of Arrays over Array of Structs}

The conventional cache entry is represented as an Array of Structs (AoS):

\begin{lstlisting}[language=Go]
// AoS: one struct per entry
type Entry struct {
    method  string     // 16 bytes (ptr + len)
    path    string     // 16 bytes
    handler uintptr    // 8 bytes
    params  []Param    // 24 bytes
}
// In memory: [entry0 | entry1 | entry2 | ...]
\end{lstlisting}

Consider a \texttt{Get} operation that must validate $n$ candidate entries by comparing their \texttt{method} and \texttt{path} fields. In AoS, each entry occupies $\sim$64+ bytes. Accessing the \texttt{method} field of entry $i$ and entry $i+1$ requires fetching two separate 64-byte cache lines. For $n$ entries, the cache-line cost is $n$ fetches just for a single field.

The Structure of Arrays (SoA) alternative, widely employed in high-performance and SIMD-oriented computing~\cite{soasimd2009}, separates fields into parallel arrays:

\begin{lstlisting}[language=Go]
// SoA: one array per field
methods  []string    // [m0, m1, m2, ...] -- contiguous
paths    []string    // [p0, p1, p2, ...] -- contiguous
handlers []uintptr   // [h0, h1, h2, ...] -- contiguous
\end{lstlisting}

When a \texttt{Get} operation scans \texttt{methods[i]} and \texttt{methods[i+1]}, both values reside in the same 64-byte cache line (a \texttt{string} header is 16 bytes, placing 4 per cache line). Accessing 4 consecutive methods costs 1 cache-line fetch instead of 4.

Formally, let $W_f$ be the width of field $f$ and $L = 64$ bytes be the cache line size. For AoS, the number of cache lines fetched to access field $f$ across $n$ entries is:

\begin{equation}
\text{Lines}_{AoS}(n, f) = n \cdot \left\lceil \frac{\text{sizeof}(\text{Entry})}{L} \right\rceil
\end{equation}

For SoA:

\begin{equation}
\text{Lines}_{SoA}(n, f) = \left\lceil \frac{n \cdot W_f}{L} \right\rceil
\end{equation}

Since $W_f \ll \text{sizeof}(\text{Entry})$ for any individual field, we have $\text{Lines}_{SoA}(n, f) \ll \text{Lines}_{AoS}(n, f)$ for all $n > 1$.

\section{Why Bitsets for Recency Tracking}

\subsection{The Cost of Pointer-Based Recency}

Traditional LRU~\cite{lru1965} tracks recency by maintaining a doubly-linked list ordered from Most Recently Used (MRU) to Least Recently Used (LRU). A cache hit requires:
\begin{enumerate}
    \item Unlinking the accessed node from its current position: 4 pointer writes.
    \item Relinking it at the MRU head: 4 pointer writes.
\end{enumerate}

Each of these pointer writes on a concurrent system must be protected by a mutual exclusion lock, serializing all \texttt{Get} operations system-wide. Lock-free linked-list variants exist but require hazard pointers or epoch-based reclamation to prevent use-after-free races~\cite{lockfree2004}, reintroducing serialization on shared epoch counters.

\subsection{CLOCK as an Approximation}

The CLOCK algorithm~\cite{clock1969,clockpro2005} approximates LRU by abandoning exact recency ordering. Instead of tracking \textit{which} entry is least recently used, it tracks \textit{whether} each entry has been accessed since the last eviction sweep---a single bit per entry.

\begin{itemize}
    \item Upon \texttt{Get}: the entry's reference bit is set to 1 via an atomic OR.
    \item Upon eviction: a hand sweeps the array. Reference bit 1 is cleared to 0 (second chance). Reference bit 0 designates the eviction victim.
\end{itemize}

This approximates LRU behavior without pointer manipulation, requiring only non-blocking atomic bitwise operations. Recent research such as S3-FIFO~\cite{s3fifo2023} and SIEVE~\cite{sieve2024} further refines the eviction policy. However, these works do not address the parallel execution cost of the eviction mechanism itself, which is the problem we solve.

\subsection{From One Bit to One 64-bit Integer}

Standard CLOCK still maintains a global clock-hand pointer---an atomic integer indicating the current sweep position. Every eviction increments this counter and examines a single entry. Under high concurrency, all evicting threads contend on this single atomic, reintroducing the $\mathcal{O}(N)$ RFO serialization we wish to avoid.

The key observation: if we partition the cache into chunks of exactly 64 entries, the valid state and access state of all 64 entries in a chunk fit into two 64-bit integers:

\begin{itemize}
    \item $V_k \in \{0,1\}^{64}$: bit $i = 1$ iff slot $i$ holds a valid entry.
    \item $A_k \in \{0,1\}^{64}$: bit $i = 1$ iff slot $i$ has been accessed since the last sweep.
\end{itemize}

A slot is an eviction candidate if it is empty ($V_k[i] = 0$) or valid but unaccessed ($V_k[i] = 1 \land A_k[i] = 0$). The complete candidate set across all 64 slots is computed in a single bitwise expression:

\begin{equation}
C_k = \neg V_k \;\lor\; (V_k \;\land\; \neg A_k)
\end{equation}

The index of the first candidate is extracted using the hardware Count Trailing Zeros (CTZ) instruction~\cite{inteltzcnt,armclz}---a single-cycle primitive on both x86\_64 (\texttt{TZCNT}) and ARM64 (\texttt{RBIT} + \texttt{CLZ}):

\begin{equation}
i^* = \text{CTZ}(C_k)
\end{equation}

This replaces an $\mathcal{O}(N)$ sequential scan over $N$ atomic loads with two register-width bitwise operations and one hardware intrinsic, bounding eviction search to a strict hardware-constant time within any chunk.

\subsection{Eliminating the Global Clock Hand}

Because \texttt{liteLRU} seeds the starting chunk index from the hash of the incoming route, $k_\text{start} = h(r) \bmod N_k$, there is no shared global clock-hand atomic. Each writer begins its eviction sweep at a deterministic but distributed starting position, naturally load-balancing eviction pressure across chunks without any coordination.

\section{Sequence Lock Design for Wait-Free Reads}

We require that \texttt{Get} operations never block on an ongoing \texttt{Add}, but also never return torn data if an eviction overwrites a slot mid-read. Sequence Locks (Seqlocks)~\cite{seqlockoriginal} achieve this. Each slot $i$ maintains an atomic 32-bit sequence counter $S_i$, initialized to 0.

\textbf{Write Protocol:} A writer sets $S_i \leftarrow S_i + 1$ (making it odd). After completing the write, it sets $S_i \leftarrow S_i + 1$ (returning to even).

\textbf{Read Protocol:} A reader samples $s_1 \leftarrow S_i$. If $s_1$ is odd, the slot is being written; report a miss immediately. Read the data. Sample $s_2 \leftarrow S_i$. If $s_1 \neq s_2$, a concurrent write occurred; report a miss. Otherwise, the read is consistent.

\textbf{Theorem (Wait-Free Reads).} \textit{The \texttt{Get} operation is wait-free: it completes in a bounded number of steps regardless of concurrent writer behavior.}

\textit{Proof.} The reader executes no loops, no spin-waits, and acquires no locks. Each decision point (odd seqlock, changed seqlock, mismatched key) results in an immediate return. The total steps are bounded by a constant: two atomic loads, one key comparison, one bitmask check. $\square$

\section{Cache-Line Padding to Eliminate False Sharing}

A naive implementation allocates seqlocks as a contiguous slice \texttt{[]atomic.Uint32}. Each \texttt{uint32} is 4 bytes, placing 16 seqlocks per 64-byte cache line. By the MESI coherency protocol~\cite{mesi1984}, a write to slot $j$ invalidates the cache line holding $S_{j+1}, \ldots, S_{j+15}$ across all other cores, inducing spurious RFOs even for logically independent operations. This is the false sharing pathology described by Bolosky and Scott~\cite{falsesharing1993}.

\textbf{Theorem (Cache Line Isolation).} \textit{If each $S_i$ occupies exactly one 64-byte cache line, concurrent operations on any two distinct slots $i \neq j$ produce no MESI coherency traffic between the cache lines of $S_i$ and $S_j$.}

\textit{Proof.} Define the padded slot state:

\begin{lstlisting}[language=Go]
type slotState struct {
    seq atomic.Uint32 //  4 bytes
    _   [60]byte      // 60 bytes padding
}                     // sizeof = 64 bytes exactly
\end{lstlisting}

Since $\text{addr}(S_{i+1}) = \text{addr}(S_i) + 64$, adjacent entries reside on disjoint, 64-byte aligned cache lines. A write to $S_i$ sets the MESI state of line $\lfloor \text{addr}(S_i) / 64 \rfloor$ to Modified, which is disjoint from line $\lfloor \text{addr}(S_j) / 64 \rfloor$ for any $j \neq i$. False sharing is structurally impossible. $\square$

The same principle is applied to the statistics counters. Rather than two globally shared \texttt{atomic.Int64} hit and miss counters, we maintain 64 independent \texttt{statStripe} structs, each padded to 64 bytes, with writers selecting a stripe via \texttt{hash \& 63}. Each core writes exclusively to its own stripe without generating coherency traffic for any other core.

\section{The Measurement Artifact of \texttt{b.RunParallel}}

Standard Go micro-benchmarks use \texttt{b.RunParallel}, distributing $b.N$ iterations across goroutines via an internal \texttt{pb.Next()} call---an atomic fetch-and-add on a shared 64-bit counter. When the operation under test completes in sub-10~ns---as is typical for a cache hit---the cost of this benchmark coordination atomic dominates the measurement.

This is a direct consequence of Amdahl's Law~\cite{amdahl1967}. Let $f$ be the sequential fraction of the workload and $N$ the number of cores. The maximum parallel speedup is:

\begin{equation}
S(N) = \frac{1}{\,(1-f) + \dfrac{f}{N}\,}
\end{equation}

When the operation time $T_\text{op}$ is comparable to the atomic latency $T_\text{atomic}$ ($\sim$20--40~ns including RFO round-trip), the \texttt{pb.Next()} counter constitutes $f \to 1$, forcing $S(N) \to 1$ regardless of $N$. The result is an apparent \textit{decrease} in throughput with additional cores---a benchmark artifact, not a property of the algorithm.

Our custom latency evaluation bypasses \texttt{pb.Next()} entirely. Each goroutine operates an independent tight loop over a pre-allocated, fixed-size key array with no shared iteration counter, eliminating the artifact and isolating the cache mechanics.

\section{Evaluation}

\textbf{Platform:} Apple M2, 8-core ARM64, macOS. Go 1.25.7. Cache capacity 1024, maxParams 20.

\textbf{Workload:} 8 concurrent workers, $1.6 \times 10^6$ total operations, 80\% \texttt{Get} / 20\% \texttt{Add}, paths drawn from a pool of 1000 distinct routes yielding $\approx$70\% hit ratio.

\subsection{Throughput Scaling}

\begin{table}[h]
\centering
\begin{tabular}{@{}rrc@{}}
\toprule
\textbf{Cores} & \textbf{Ops/sec} & \textbf{Scaling Factor} \\ \midrule
1              & 5,457,508        & 1.00$\times$            \\
2              & 8,847,292        & 1.62$\times$            \\
4              & 13,701,548       & 2.51$\times$            \\
8              & 17,066,659       & 3.13$\times$            \\ \bottomrule
\end{tabular}
\caption{Throughput under un-synchronized worker methodology.}
\end{table}

Sub-linear scaling from 4 to 8 cores is expected: at 8 cores on an M2 die, the shared L3 and memory bus bandwidth begin to constrain throughput independently of lock contention.

\subsection{Latency Percentiles}

\begin{table}[h]
\centering
\begin{tabular}{@{}ccc@{}}
\toprule
\textbf{Percentile} & \textbf{Measured} & \textbf{Est. (40~ns overhead removed)} \\ \midrule
p50                 & 250~ns            & $\sim$210~ns                           \\
p99                 & 1{,}000~ns        & $\sim$960~ns                           \\
p99.9               & 1{,}417~ns        & $\sim$1{,}377~ns                       \\
Max                 & $\sim$7.1~ms      & ---                                    \\ \bottomrule
\end{tabular}
\caption{Latency percentiles under 8-core contended workload. Measurement overhead due to \texttt{time.Now()} is $\sim$40~ns per sample.}
\end{table}

The observed maximum latency is dominated by OS goroutine scheduling jitter, not cache mechanics.

\section{Discussion and Limitations}

\textbf{Approximate Eviction.} The Chunked Bitmask CLOCK algorithm is an approximation of LRU~\cite{lru1965,clock1969}. In workloads with access patterns adversarially tuned to the CLOCK approximation error, hit-rate may degrade relative to true LRU or more recent policies such as SIEVE~\cite{sieve2024}. The tradeoff is deliberate: structural lock freedom is valued over exact eviction fidelity for the target workload of web-scale routing caches.

\textbf{Fixed Capacity.} The cache does not resize. The hash map is pre-allocated at construction. This prevents any allocation or growth during steady-state execution but requires the caller to correctly provision capacity at initialization time.

\textbf{String Copying.} While \texttt{Get} performs zero heap allocations for the routing lookup, returning a \texttt{[]Param} copy requires a pool-managed allocation. The sync.Pool bounds this cost but it is not zero in all cases.

\section{Conclusion}

We have presented the hardware-grounded derivation of each principal design decision in \texttt{liteLRU}: off-heap mmap allocation to eliminate GC interference~\cite{gogc2018}; SoA layout to maximize cache-line occupancy per field access; bitset-encoded recency state enabling bulk $\mathcal{O}(1)$ evaluation via hardware \texttt{CTZ} intrinsics~\cite{inteltzcnt,armclz}; Seqlocks~\cite{seqlockoriginal} to provide wait-free read semantics; and 64-byte cache-line padding~\cite{falsesharing1993,mesi1984} to eliminate false-sharing serialization. The result is a cache architecture whose concurrency cost is bounded by single-cycle hardware instructions rather than operating system synchronization primitives.

\bibliographystyle{plain}
\bibliography{references}

\end{document}
